% !TeX root = ../../tesis.tex

% =============================================================================
%  B.3 — CONSTRUCTOR DEL GRAFO: VFTGraphBuilder
% =============================================================================
\section{Constructor del Grafo Topológico (\texttt{VFTGraphBuilder})}
\label{anx:vft-graphbuilder}

\texttt{VFTGraphBuilder} es la clase que materializa el grafo dirigido
(\texttt{nx.DiGraph}) a partir del \textsc{geojson} ya validado por Pydantic.
La construcción ocurre en tres fases secuenciales y una opcional:

\begin{enumerate}
    \item \textbf{Fase 1 — Nodos:} cada \texttt{Feature} de tipo
    \texttt{estacion} se registra como nodo con sus coordenadas GPS y
    atributos operativos.
    \item \textbf{Fase 2 — Aristas de tránsito:} un \termino{KDTree} por
    sistema recorre cada \texttt{MultiLineString} y crea una arista entre
    cada par de estaciones consecutivas, acumulando la distancia Haversine
    real del trazo.
    \item \textbf{Filtro fantasma:} elimina aristas de tránsito con
    distancia superior a 5\,km, causadas por geometrías
    \texttt{MultiLineString} fragmentadas en el servidor Go.
    \item \textbf{Fase 3 — Transbordos peatonales} (solo en modo
    \texttt{REALISTIC\_INTEGRATION}): aristas de caminata entre estaciones
    de distintos sistemas dentro de la tolerancia configurada, cobrando
    tiempo de espera al siguiente servicio (\termino{boarding cost}).
\end{enumerate}

\subsection*{Umbrales estadísticos de transbordo}

Los umbrales pre-configurados se derivan de la distribución empírica de
distancias entre estaciones del \textsc{geojson} de Apimetro
(corrida 2026-04-18). El umbral por defecto para análisis es $Q_1 = 85$\,m,
que representa el transbordo directo más conservador.

\begin{table}[H]
    \centering
    \small
    \caption[Umbrales estadísticos de transbordo peatonal]{Umbrales estadísticos
    de transbordo peatonal en VFTModel.
    \fuentefigura{Elaboración propia con base en \texttt{graph\_builder.py}
    (VFTModel, 2026).}}
    \label{tab:anx-vft-umbrales}
    \begin{tabular}{|l|r|l|}
        \hline
        \textbf{Clave} & \textbf{Valor (m)} & \textbf{Interpretación} \\
        \hline
        \texttt{MIN}    &   15 & Intersección estricta / mismo polígono arquitectónico \\
        \texttt{Q1}     &   85 & Transbordo rápido y directo (conservador) \\
        \texttt{Q2}     &  180 & Transbordo promedio real — integración plena de CETRAMs \\
        \texttt{MEAN}   &  245 & Promedio matemático (sesgado por valores extremos) \\
        \texttt{Q3}     &  420 & Límite máximo de caminata forzada \\
        \texttt{MAX}    &  880 & Valor extremo / falso transbordo \\
        \hline
    \end{tabular}
\end{table}

\subsection*{Construcción de la red base (Fases 1 y 2)}

\begin{lstlisting}[style=estiloPython,
    caption={Construcción de nodos y aristas de tránsito mediante KDTree por
             sistema. \texttt{src/core/services/graph\_builder.py}.},
    label=lst:vft-graph-base]
# Fase 1: registrar estaciones como nodos
for feature in self.validated_data.features:
    if feature.properties.tipo_entidad == TipoEntidad.estacion:
        lon, lat   = feature.geometry.coordinates
        node_id    = (lon, lat)
        self.G.add_node(node_id,
            pos=(lon, lat),
            nombre=feature.properties.nombre,
            sistema=feature.properties.sistema.value,
            es_cetram=feature.properties.es_cetram,
            alcaldia_municipio=feature.properties.alcaldia_municipio)

# Fase 2: KDTree por sistema — evita contaminación cruzada entre redes
sys_kdtrees = {
    s: (KDTree(np.array([[n[0], n[1]] for n in nlist])), nlist)
    for s, nlist in sys_nodes_map.items() if nlist
}
# Para cada ruta: acumular distancia haversine entre waypoints y
# crear aristas entre estaciones consecutivas detectadas por KDTree
for i, coord in enumerate(all_coords):
    dist_acumulada += haversine(all_coords[i-1], coord)  # tramo real
    dist_deg, idx   = kdtree_s.query([lon, lat])
    if dist_deg <= SNAP_TOLERANCE_DEG:  # ~50 m
        if estacion_node != ultimo_waypoint:
            self.G.add_edge(ultimo_waypoint, estacion_node,
                distancia_segmento_m=round(dist_acumulada, 2), **edge_attr)
\end{lstlisting}

\subsection*{Transbordos peatonales con \termino{boarding cost} (Fase 3)}

\begin{lstlisting}[style=estiloPython,
    caption={Cálculo del peso de aristas peatonales: caminata + espera al
             siguiente servicio. \texttt{src/core/services/graph\_builder.py}.},
    label=lst:vft-snapping]
for i in range(len(nodos_estacion)):
    for j in range(i + 1, len(nodos_estacion)):
        distancia_m = haversine(u_attr['pos'], v_attr['pos'])
        if 0 < distancia_m <= tolerance_m:
            tiempo_caminata = distancia_m / (5000.0 / 60.0)  # 5 km/h
            wait_v = FALLBACK_FRECUENCIA.get(sistema_v, 10.0) / 2.0
            wait_u = FALLBACK_FRECUENCIA.get(sistema_u, 10.0) / 2.0
            # Arista ida: caminar hacia V + esperar servicio V
            self.G.add_edge(u_id, v_id, tipo="transfer",
                weight=round(tiempo_caminata + wait_v, 4))
            # Arista vuelta: caminar hacia U + esperar servicio U
            self.G.add_edge(v_id, u_id, tipo="transfer",
                weight=round(tiempo_caminata + wait_u, 4))
\end{lstlisting}
